% Section 20 Header
\section{Pattern 18: Advanced Patterns}

\begin{frame}[c]{\empty}
    \begin{center}
        \Large\textbf{\secname}
    \end{center}
\end{frame}

% Slide 1: Concept
\begin{frame}{Advanced Algorithmic Patterns}
  For Senior and Staff level roles, interviewers often present advanced structures:
  \begin{itemize}
    \item \textbf{Segment Tree:} A binary tree used for storing intervals or segments. It allows querying a range (sum, min, max) and modifying elements in $O(\log N)$ time.
    \item \textbf{Topological Sort:} Ordering vertices in a directed acyclic graph (DAG) such that for every directed edge $uv$, vertex $u$ comes before $v$. (Kahn's BFS or DFS).
    \item \textbf{Sweep Line:} A geometry paradigm that solves range/interval problems by sorting boundary events (start/end) and sweeping a imaginary line across them.
    \item \textbf{Trie (Prefix Tree):} A search tree used to store associative keys (typically strings) efficiently. Excellent for autocompletion, spelling checkers, prefix matching.
  \end{itemize}
\end{frame}

% Problem 1: Segment Tree
\begin{frame}[c]{\empty}
    \begin{center}
        \Large \textbf{Segment Tree: \href{https://leetcode.com/problems/range-sum-query-mutable/}{Range Sum Query - Mutable}}
    \end{center}
\end{frame}

% Problem 2: Topological Sort
\begin{frame}[c]{\empty}
    \begin{center}
        \Large \textbf{Topological Sort: \href{https://leetcode.com/problems/course-schedule/}{Course Schedule} \& \href{https://leetcode.com/problems/course-schedule-ii/}{Course Schedule II}}
    \end{center}
\end{frame}

% Problem 3: Sweep Line
\begin{frame}[c]{\empty}
    \begin{center}
        \Large \textbf{Sweep Line: \href{https://leetcode.com/problems/the-skyline-problem/}{The Skyline Problem}}
    \end{center}
\end{frame}

% Problem 4: Tries
\begin{frame}[c]{\empty}
    \begin{center}
        \Large \textbf{Tries: \href{https://leetcode.com/problems/implement-trie-prefix-tree/}{Implement Trie} \& \href{https://leetcode.com/problems/replace-words/}{Replace Words}}
    \end{center}
\end{frame}
